% =============================================================================
% Abstract — conservative, evidence-aligned framing
% Drop-in include: % =============================================================================
% Abstract — conservative, evidence-aligned framing
% Drop-in include: % =============================================================================
% Abstract — conservative, evidence-aligned framing
% Drop-in include: % =============================================================================
% Abstract — conservative, evidence-aligned framing
% Drop-in include: \input{sections/abstract}
% Reviewed-record descendant; target: 150--180 words
% =============================================================================
\begin{abstract}
Group-relative RL can lose its reward-driven learning signal when every
completion for a prompt receives the same reward. This is a diagnostic audit and
correction record, not a new algorithm, benchmark, or algorithm ranking. It is a
descendant of NeurIPS submission \#36320, not its authenticated or byte-identical
reviewed PDF; we audit that submitted record rather than claiming that local edits
change it (overlap map in Section~\ref{sec:intro}). The descriptive reporting field
Zero-Variance Fraction (ZVF) records homogeneous reward groups; its complement
$1-\mathrm{ZVF}$ is the mixed-group fraction (historical alias: Gradient
Utilization, GU, which measures neither gradient magnitude nor usefulness). Whether
reward contrast, rather than algorithm labels, explains cross-stack gaps is a
hypothesis this audit does \emph{not} demonstrate: no same-stack identified contrast
manipulation is reported.
In the two reviewed synthetic tool-use cells, reward is 0, ZVF is 1, GU is 0,
and held-out performance was not evaluated. These are two descriptive all-wrong
cases, not a validated predictor or stopping policy. For the submitted
group-standardized runner, homogeneous rewards make the centered reward-contrast
term zero. The runner is not canonical GRPO: it omits ratio/clipping, reference-KL,
and completion-only masking, so this identity does not imply a zero total gradient
or transfer to other implementations. GSM8K and later per-run artifacts are kept
separate by estimand and provenance; the historical heterogeneous checkpoint mean
is withdrawn. We provide an evidence-tiered reporting audit and a prospective
operator-decision contract, not a deployment benefit or algorithm-ranking claim.
\end{abstract}

% Reviewed-record descendant; target: 150--180 words
% =============================================================================
\begin{abstract}
Group-relative RL can lose its reward-driven learning signal when every
completion for a prompt receives the same reward. This is a diagnostic audit and
correction record, not a new algorithm, benchmark, or algorithm ranking. It is a
descendant of NeurIPS submission \#36320, not its authenticated or byte-identical
reviewed PDF; we audit that submitted record rather than claiming that local edits
change it (overlap map in Section~\ref{sec:intro}). The descriptive reporting field
Zero-Variance Fraction (ZVF) records homogeneous reward groups; its complement
$1-\mathrm{ZVF}$ is the mixed-group fraction (historical alias: Gradient
Utilization, GU, which measures neither gradient magnitude nor usefulness). Whether
reward contrast, rather than algorithm labels, explains cross-stack gaps is a
hypothesis this audit does \emph{not} demonstrate: no same-stack identified contrast
manipulation is reported.
In the two reviewed synthetic tool-use cells, reward is 0, ZVF is 1, GU is 0,
and held-out performance was not evaluated. These are two descriptive all-wrong
cases, not a validated predictor or stopping policy. For the submitted
group-standardized runner, homogeneous rewards make the centered reward-contrast
term zero. The runner is not canonical GRPO: it omits ratio/clipping, reference-KL,
and completion-only masking, so this identity does not imply a zero total gradient
or transfer to other implementations. GSM8K and later per-run artifacts are kept
separate by estimand and provenance; the historical heterogeneous checkpoint mean
is withdrawn. We provide an evidence-tiered reporting audit and a prospective
operator-decision contract, not a deployment benefit or algorithm-ranking claim.
\end{abstract}

% Reviewed-record descendant; target: 150--180 words
% =============================================================================
\begin{abstract}
Group-relative RL can lose its reward-driven learning signal when every
completion for a prompt receives the same reward. This is a diagnostic audit and
correction record, not a new algorithm, benchmark, or algorithm ranking. It is a
descendant of NeurIPS submission \#36320, not its authenticated or byte-identical
reviewed PDF; we audit that submitted record rather than claiming that local edits
change it (overlap map in Section~\ref{sec:intro}). The descriptive reporting field
Zero-Variance Fraction (ZVF) records homogeneous reward groups; its complement
$1-\mathrm{ZVF}$ is the mixed-group fraction (historical alias: Gradient
Utilization, GU, which measures neither gradient magnitude nor usefulness). Whether
reward contrast, rather than algorithm labels, explains cross-stack gaps is a
hypothesis this audit does \emph{not} demonstrate: no same-stack identified contrast
manipulation is reported.
In the two reviewed synthetic tool-use cells, reward is 0, ZVF is 1, GU is 0,
and held-out performance was not evaluated. These are two descriptive all-wrong
cases, not a validated predictor or stopping policy. For the submitted
group-standardized runner, homogeneous rewards make the centered reward-contrast
term zero. The runner is not canonical GRPO: it omits ratio/clipping, reference-KL,
and completion-only masking, so this identity does not imply a zero total gradient
or transfer to other implementations. GSM8K and later per-run artifacts are kept
separate by estimand and provenance; the historical heterogeneous checkpoint mean
is withdrawn. We provide an evidence-tiered reporting audit and a prospective
operator-decision contract, not a deployment benefit or algorithm-ranking claim.
\end{abstract}

% Reviewed-record descendant; target: 150--180 words
% =============================================================================
\begin{abstract}
Group-relative RL can lose its reward-driven learning signal when every
completion for a prompt receives the same reward. This is a diagnostic audit and
correction record, not a new algorithm, benchmark, or algorithm ranking. It is a
descendant of NeurIPS submission \#36320, not its authenticated or byte-identical
reviewed PDF; we audit that submitted record rather than claiming that local edits
change it (overlap map in Section~\ref{sec:intro}). The descriptive reporting field
Zero-Variance Fraction (ZVF) records homogeneous reward groups; its complement
$1-\mathrm{ZVF}$ is the mixed-group fraction (historical alias: Gradient
Utilization, GU, which measures neither gradient magnitude nor usefulness). Whether
reward contrast, rather than algorithm labels, explains cross-stack gaps is a
hypothesis this audit does \emph{not} demonstrate: no same-stack identified contrast
manipulation is reported.
In the two reviewed synthetic tool-use cells, reward is 0, ZVF is 1, GU is 0,
and held-out performance was not evaluated. These are two descriptive all-wrong
cases, not a validated predictor or stopping policy. For the submitted
group-standardized runner, homogeneous rewards make the centered reward-contrast
term zero. The runner is not canonical GRPO: it omits ratio/clipping, reference-KL,
and completion-only masking, so this identity does not imply a zero total gradient
or transfer to other implementations. GSM8K and later per-run artifacts are kept
separate by estimand and provenance; the historical heterogeneous checkpoint mean
is withdrawn. We provide an evidence-tiered reporting audit and a prospective
operator-decision contract, not a deployment benefit or algorithm-ranking claim.
\end{abstract}
